% =====================================================================
% Jacobi + Newton para trazas de matrices, con una aplicacion a Ihara.
% Companero de las formalizaciones Godsil--Gutman (Paper I), Heilmann--Lieb
% (Paper II) e Ihara--Bass. Mismo estilo de la casa: tecnico, inequivoco,
% guiado por una historia (cada seccion abre con una pregunta), formulas clave
% enmarcadas y centradas, humildad en todo. Reglas de honestidad literales:
% "demostrado" = sin sorry; alcance exacto; nada oculto.
% Verificado contra la fuente Lean 2026-06-07 (lake build exit 0; #print axioms
% reporta propext, Classical.choice, Quot.sound). Toolchain v4.30.0-rc2.
% =====================================================================
\documentclass[11pt]{article}
\usepackage[T1]{fontenc}
\usepackage[spanish,es-noindentfirst]{babel}
\usepackage{lmodern}
\usepackage{amsmath,amssymb,amsthm,booktabs,graphicx,hyperref}
\usepackage{xurl}
\usepackage[table]{xcolor}
\definecolor{ggteal}{HTML}{1B6F8C}
\definecolor{ggband}{HTML}{DCE7EB}
\definecolor{ggamber}{HTML}{E08A1E}
\usepackage[margin=1in]{geometry}
\usepackage{microtype}
\usepackage[most]{tcolorbox}
\usepackage{tikz}
\usetikzlibrary{arrows.meta,positioning,calc}
\setlength{\emergencystretch}{2em}
\newtheorem{theorem}{Teorema}
\newtheorem{lemma}{Lema}
\newtheorem{definition}{Definici\'on}
\newcommand{\lean}[1]{\texttt{#1}}
\newcommand{\Z}{\mathbb{Z}}
\newcommand{\N}{\mathbb{N}}
\newcommand{\tr}{\operatorname{tr}}
\newcommand{\adj}{\operatorname{adj}}
\DeclareMathOperator{\charpolyRev}{charpolyRev}
\hypersetup{colorlinks=true,linkcolor=ggteal,citecolor=ggteal,urlcolor=ggteal}
\newtcbox{\keybox}{on line,boxrule=0.9pt,colframe=ggteal,colback=ggband!40,
  arc=3pt,left=8pt,right=8pt,top=5pt,bottom=5pt}
\newcommand{\keyeq}[1]{\begin{center}\keybox{$\displaystyle #1$}\end{center}}

\title{\bf Lo que sabe la derivada de un determinante\\[4pt]
  \large\normalfont La f\'ormula de Jacobi y la identidad de Newton para trazas de matrices
  en Lean~4, con una aplicaci\'on a los conteos de Ihara}
\author{Carles Mar\'in\thanks{Formalizado con asistencia de IA (Claude, Anthropic); las matem\'aticas
y todas las afirmaciones son responsabilidad del autor. La metodolog\'ia se discute en la
Secci\'on~\ref{sec:implications}.}\\[3pt]
  \normalsize Investigador independiente\\
  \normalsize \texttt{karlesmarin@gmail.com}}
\date{Junio de 2026}

\begin{document}
\maketitle

\begin{abstract}
Una regla de una sola l\'inea de un curso de \'algebra lineal---\emph{la derivada de un determinante
es la traza de la adjugada por la derivada de la matriz}---resulta ser la semilla de la que crecen
los coeficientes del polinomio caracter\'istico, las sumas de potencias de las trazas de una matriz
y, en \'ultima instancia, los conteos de caminos cerrados de un grafo. Este art\'iculo formaliza esa
semilla y sus primeros frutos en Lean~4 / Mathlib, todo sin \lean{sorry}. El resultado principal es
la \emph{f\'ormula de Jacobi} $(\det M)'=\tr(\adj M\cdot M')$ para una matriz de polinomios, de la
que no he hallado demostraci\'on previa en ning\'un asistente de pruebas; Mathlib s\'olo ten\'ia la
derivada \emph{anal\'itica} de $\det$ entre espacios normados, nunca la identidad algebraica. De
ella derivo la \emph{identidad de Newton para trazas de matrices}---la derivada logar\'itmica del
polinomio caracter\'istico invertido re\'une las sumas de potencias de las trazas---, un enunciado a
todos los \'ordenes que Mathlib ofrec\'ia s\'olo en el coeficiente lineal, y que no necesita
autovalores. El motor es una \emph{serie resolvente} $(I-XM)^{-1}=\sum_k M^kX^k$ sobre series formales
de potencias, tambi\'en nueva en la biblioteca. Como aplicaci\'on las compongo con mi formalizaci\'on
previa de la f\'ormula del determinante de Bass para leer los conteos no-retroceso de Ihara
$N_k=\tr(B^k)$ desde el espectro de adyacencia, y para demostrar combinatoriamente que $\tr(B^k)$
cuenta caminos cerrados sin retroceso. Soy exacto con la frontera: nada de esto es matem\'atica
nueva; el conteo de geod\'esicas primas que la estructura c\'iclica prepara s\'olo se construye hasta
el peso invariante por rotaci\'on, no hasta el producto de Euler; y $\tr(B^k)$ cuenta caminos cerrados
\emph{con punto base}, no clases geod\'esicas. La contribuci\'on es la formalizaci\'on, y una peque\~na
pieza de infraestructura certificada sobre la que cualquier desarrollo futuro de polinomios
caracter\'isticos, funciones zeta o espectros de expansores podr\'a apoyarse.
\end{abstract}

\noindent\emph{Mapa para el lector.} Si s\'olo lees dos cosas, lee el Teorema~\ref{thm:jacobi}
(Jacobi) y el Teorema~\ref{thm:newton} (Newton): el primero dice que la derivada de un determinante
es una traza; el segundo dice que, a todos los \'ordenes, ese \'unico hecho \emph{es} el puente entre
las trazas de una matriz y su polinomio caracter\'istico. Todo lo dem\'as---el motor resolvente
(Secci\'on~\ref{sec:resolvent}), el conteo de caminos (Secci\'on~\ref{sec:walks}), la aplicaci\'on a
Ihara (Secci\'on~\ref{sec:ihara})---cuelga de esos dos. Las f\'ormulas que cargan el peso van
enmarcadas; los l\'imites honestos se enuncian con la misma claridad que los resultados.

\begin{figure}[h]
\centering
\begin{tikzpicture}[
  node distance=6mm and 9mm,
  box/.style={draw=ggteal,fill=ggband!40,rounded corners=3pt,line width=0.8pt,
    align=center,inner sep=5pt,text width=27mm,font=\small},
  arr/.style={-{Stealth[length=2.5mm]},ggteal,line width=0.9pt},
  lab/.style={font=\scriptsize\itshape,ggamber,align=center}]
\node[box] (jac) {\textbf{Jacobi}\\$(\det M)'=\tr(\adj M\,M')$};
\node[box,right=of jac] (new) {\textbf{Newton}\\$\sum_k\tr(M^k)X^k$\\$=-X\,\tfrac{(\charpolyRev)'}{\charpolyRev}$};
\node[box,right=of new] (bass) {\textbf{$+$ Bass}\\$\det(I-XB)$\\$\to$ datos $(A,D)$};
\node[box,right=of bass] (nk) {\textbf{Ihara}\\$N_k=\tr(B^k)$\\$=\#$caminos};
\draw[arr] (jac) -- node[lab,above]{derivar a\\todo orden} (new);
\draw[arr] (new) -- node[lab,above]{componer} (bass);
\draw[arr] (bass) -- node[lab,above]{deriv. log.} (nk);
\end{tikzpicture}
\caption{La columna vertebral del art\'iculo. Un hecho elemental---la derivada de un determinante es
una traza (Jacobi)---, llevado a todos los \'ordenes, se vuelve el diccionario entre las trazas de una
matriz y su polinomio caracter\'istico (Newton), y compuesto con la f\'ormula del determinante de
Bass alcanza los conteos de caminos cerrados sin retroceso $N_k=\tr(B^k)$ de un grafo. Cada flecha es
un teorema de este art\'iculo.}
\label{fig:spine}
\end{figure}

% ---------------------------------------------------------------
\section{Introducci\'on: una breve historia, y por qu\'e este art\'iculo}
\label{sec:intro}

\emph{?`Podr\'ia una regla para derivar un determinante---un hecho tan elemental que suele ponerse
como ejercicio---ser el motor que, llevado a todos los \'ordenes, recupera el polinomio
caracter\'istico desde las trazas de una matriz, y luego cuenta los caminos cerrados de un grafo?} La
respuesta es s\'i, y seguirla exige que se encuentren tres breves historias.

La primera es la \emph{f\'ormula de Jacobi}. En el siglo XIX Jacobi registr\'o c\'omo cambia el
determinante de una matriz al variar sus entradas: $\mathrm{d}(\det M)=\tr(\adj M\cdot\mathrm{d}M)$,
donde $\adj M$ es la adjugada (la traspuesta de la matriz de cofactores). Es el c\'alculo matricial
tras la f\'ormula de Liouville para el wronskiano de una EDO lineal, tras el algoritmo de
Faddeev--LeVerrier que computa un polinomio caracter\'istico desde las trazas, y tras toda
perturbaci\'on a primer orden de un autovalor. Es tambi\'en, en su forma \emph{algebraica}---sobre un
anillo conmutativo arbitrario, con la derivada formal de polinomios en vez de un l\'imite---, una
ausente de la mayor biblioteca formal de matem\'aticas. Mathlib sab\'ia derivar un determinante
\emph{anal\'iticamente}, como aplicaci\'on suave entre espacios normados; no conoc\'ia la identidad
algebraica de una l\'inea.

La segunda son las \emph{identidades de Newton}. Newton relacion\'o las sumas de potencias
$p_k=\sum_i\lambda_i^k$ de una colecci\'on de n\'umeros con sus funciones sim\'etricas elementales
$e_k$---equivalentemente, los coeficientes del polinomio $\prod_i(X-\lambda_i)$. Especializadas a los
autovalores de una matriz, las sumas de potencias son exactamente las trazas $\tr(M^k)$, y las
funciones sim\'etricas elementales son exactamente los coeficientes del polinomio caracter\'istico.
As\'i, las identidades de Newton, en esta forma, son un puente: permiten computar el polinomio
caracter\'istico desde las trazas de las potencias, y rec\'iprocamente. Mathlib tiene las identidades
de Newton, pero s\'olo en abstracto---para polinomios sim\'etricos en variables libres---, nunca
conectadas con las trazas de una matriz. El puente exist\'ia; su extremo matricial faltaba, registrado
s\'olo en el primer coeficiente como $\tr M=-\,(\text{coeficiente siguiente del polinomio
caracter\'istico})$.

El tercer hilo enlaza esto con los grafos, y con mi trabajo anterior. La funci\'on zeta de Ihara de
un grafo cuenta sus caminos cerrados sin retroceso; su rec\'iproco es el determinante $\det(I-XB)$
del operador no-retroceso $B$ de Hashimoto, y formalic\'e la reducci\'on de \emph{Bass} de ese
determinante a datos de tama\~no v\'ertice en un art\'iculo compa\~nero~\cite{IharaBass}. Los conteos
$N_k=\tr(B^k)$ son exactamente las sumas de potencias de trazas que gobierna la identidad de Newton.
La misma derivada que mueve a Jacobi y a Newton alcanza, por tanto, al aplicarse a la identidad de
Bass, los conteos de ciclos de un grafo.

Este art\'iculo formaliza esa cadena---Jacobi, luego Newton, luego los conteos de Ihara---en Lean~4, y
es exacto sobre d\'onde, por ahora, termina el camino.

% ---------------------------------------------------------------
\section{La f\'ormula de Jacobi}
\label{sec:jacobi}

\emph{Si empujas una entrada de una matriz, ?`c\'omo se mueve su determinante---y por qu\'e habr\'ia
de involucrar la respuesta la traza de algo?}

Para una matriz $M$ cuyas entradas son polinomios en una variable formal, escribe $M'$ para la
derivada entrada a entrada y $\det M$ para el determinante (un polinomio). La respuesta es una traza.

\begin{theorem}[F\'ormula de Jacobi, en Lean~4]
\label{thm:jacobi}
Para una matriz $M$ de $n\times n$ de polinomios sobre un anillo conmutativo,
\keyeq{(\det M)'\;=\;\tr\!\bigl(\adj M\cdot M'\bigr).}
El enunciado es sin \lean{sorry} (\lean{Matrix.derivative\_det}); \lean{\#print axioms} reporta
\'unicamente \lean{propext}, \lean{Classical.choice}, \lean{Quot.sound}.
\end{theorem}

La demostraci\'on es la propia multilinealidad del determinante, hecha columna a columna. Derivando
$\det M=\sum_\sigma\varepsilon_\sigma\prod_i M_{\sigma(i),i}$ por la regla de Leibniz y reagrupando se
ve que $(\det M)'$ es la suma, sobre columnas $k$, del determinante de $M$ con la columna $k$
reemplazada por su derivada; la regla de Cramer lee cada uno de esos determinantes de una columna en
la adjugada, y la suma se reensambla como una traza. (En Lean el reagrupado usa la forma de producto
de la f\'ormula de Leibniz, \lean{det\_apply'}, para mantener cada coeficiente un elemento genuino del
anillo en vez de una permutaci\'on con signo---una peque\~na elecci\'on que vuelve toda la
demostraci\'on una cadena de reescrituras.)

\paragraph{Por qu\'e es \'util, llanamente.}
La derivada de un determinante aparece donde sea que una matriz se mueve. Tres lugares donde se gana
el sustento: la \emph{f\'ormula de Liouville}---el wronskiano de un sistema lineal evoluciona como
$\det(\Phi)'=\tr(A)\,\det(\Phi)$, que es la f\'ormula de Jacobi para $\Phi'=A\Phi$; la
\emph{perturbaci\'on de autovalores}---el cambio a primer orden de un autovalor simple es una traza
contra su proyecci\'on espectral; y el \emph{algoritmo de Faddeev--LeVerrier}, que extrae un polinomio
caracter\'istico y una inversa de matriz s\'olo a partir de trazas de potencias, y cuya correcci\'on
descansa exactamente en esta identidad. Cada uno es un teorema que alguien podr\'a un d\'ia querer
certificar; cada uno tiene ya su primera piedra.

% ---------------------------------------------------------------
\section{La serie resolvente: un motor sin autovalores}
\label{sec:resolvent}

\emph{?`Se puede hablar de $(I-XM)^{-1}$ sin invertir nunca nada---y sin mencionar jam\'as un
autovalor?}

Sobre series formales de potencias, s\'i. La serie geom\'etrica $\sum_k M^kX^k$ es una matriz honesta
de series de potencias, y es la inversa de $I-XM$.

\begin{definition}[Serie resolvente]
Para una matriz $M$ de $n\times n$ sobre un anillo conmutativo, la \emph{serie resolvente} es la
matriz de series formales de potencias $\bigl(\,\sum_k M^k X^k\,\bigr)$, escrita
\lean{resolventSeries M}.
\end{definition}

\begin{lemma}[La resolvente invierte $I-XM$, y su traza es la funci\'on generatriz de trazas]
\label{lem:resolvent}
\keyeq{(I-XM)\cdot\textstyle\sum_k M^kX^k\;=\;I,\qquad
   \tr\!\bigl(\textstyle\sum_k M^kX^k\bigr)\;=\;\sum_k \tr(M^k)\,X^k.}
Ambas sin \lean{sorry} (\lean{one\_sub\_smul\_mul\_resolventSeries}, \lean{trace\_resolventSeries}). La
identidad de la inversa se demuestra a partir de la ecuaci\'on de punto fijo $R=I+XM\!\cdot\!R$ por un
\'unico reordenamiento---sin autovalores, sin clausura algebraica, sobre cualquier anillo conmutativo.
\end{lemma}

Este es el motor de la Secci\'on~\ref{sec:newton}. Su gracia es justamente su modestia: la forma
est\'andar de relacionar las trazas de potencias con el polinomio caracter\'istico pasa por los
autovalores, y as\'i necesita que la matriz se descomponga---un cuerpo algebraicamente cerrado, o un
rodeo por la matriz gen\'erica. La serie resolvente prescinde de todo ello. La traza de $(I-XM)^{-1}$
es, por construcci\'on, la \emph{funci\'on generatriz de trazas} $\sum_k\tr(M^k)X^k$, y ese es el
objeto del que trata la identidad de Newton.

\paragraph{Por qu\'e es \'util, llanamente.}
Una resolvente formal es el hogar natural de cualquier cantidad que uno quiera leer ``orden a orden''
en una perturbaci\'on: los momentos $\tr(M^k)$ de una matriz, la funci\'on de Green de un operador
discreto, la funci\'on generatriz de caminos (Secci\'on~\ref{sec:walks}). Tener $(I-XM)^{-1}$
disponible como infraestructura certificada de series de potencias---en lugar de rederivar la serie
geom\'etrica cada vez---es una peque\~na comodidad que se acumula.

% ---------------------------------------------------------------
\section{La identidad de Newton para trazas de matrices}
\label{sec:newton}

\emph{Las trazas $\tr(M),\tr(M^2),\dots$ y los coeficientes del polinomio caracter\'istico son dos
descripciones de la misma matriz. ?`Cu\'al es, exactamente, el diccionario entre ellas?}

Es una sola identidad de series de potencias, y es la f\'ormula de Jacobi llevada a todos los
\'ordenes. Escribe $\charpolyRev M$ para el polinomio caracter\'istico invertido $\det(I-XM)$---un
polinomio con t\'ermino constante $1$, y por tanto una unidad en el anillo de series.

\begin{theorem}[Identidad de Newton para trazas de matrices, en Lean~4]
\label{thm:newton}
Para una matriz $M$ de $n\times n$ sobre un anillo conmutativo, en el anillo de series formales de
potencias,
\keyeq{(\charpolyRev M)'\;=\;-\,\charpolyRev M\cdot\sum_{k\ge0}\tr(M^{k+1})\,X^k.}
Equivalentemente $\sum_{k\ge1}\tr(M^k)X^k=-X\,(\charpolyRev M)'/\charpolyRev M$, la derivada
logar\'itmica. El enunciado es sin \lean{sorry} (\lean{Matrix.charpolyRev\_logDeriv}), sobre cualquier
anillo conmutativo, sin autovalores. En el coeficiente lineal se especializa al lema que Mathlib ya
ten\'ia, $\tr M=-\,[\,X^1\,]\,\charpolyRev M$.
\end{theorem}

La demostraci\'on es la cadena prometida en la introducci\'on. Deriva $\charpolyRev M=\det(I-XM)$ por
Jacobi (Teorema~\ref{thm:jacobi}); la derivada de $I-XM$ es $-M$, as\'i que $(\charpolyRev M)'$ es
$-\tr(\adj(I-XM)\cdot M)$. La adjugada de $I-XM$ es $\det(I-XM)$ por su inversa---que, por el
Lema~\ref{lem:resolvent}, es $\charpolyRev M$ por la serie resolvente; la traza de la resolvente
contra $M$ es la funci\'on generatriz desplazada $\sum_k\tr(M^{k+1})X^k$. Transportar la identidad
polinomial de Jacobi por la coerci\'on hacia las series de potencias y sustituir cierra el lazo. Todo
el argumento nunca nombra un autovalor.

\paragraph{Por qu\'e es \'util, llanamente.}
Este es el diccionario mismo. Le\'ido en un sentido, computa los coeficientes del polinomio
caracter\'istico---y por tanto el determinante, la traza y (v\'ia Cayley--Hamilton) la
inversa---desde las trazas de las potencias; le\'ido en el otro, computa esas trazas desde el
polinomio. Es la forma matricial de las identidades de Newton, la forma que un matem\'atico encuentra
de verdad, y la forma que la versi\'on de funciones sim\'etricas de Mathlib no alcanzaba. La forma de
derivada logar\'itmica $\sum\tr(M^k)X^k=-X(\charpolyRev)'/\charpolyRev$ es, adem\'as, el patr\'on
universal tras las \emph{funciones zeta}: el rec\'iproco de un determinante, derivado
logar\'itmicamente, da un conteo---que es precisamente como la Secci\'on~\ref{sec:ihara} alcanza el
grafo.

% ---------------------------------------------------------------
\section{Las potencias de matrices cuentan caminos}
\label{sec:walks}

\emph{Antes de contar ciclos, ?`qu\'e cuenta una sola entrada de $M^k$?}

La respuesta cl\'asica---caminos---necesita un hogar formal, y para grafos \emph{dirigidos} Mathlib no
lo ten\'ia. Registra que una potencia de la matriz de adyacencia de un grafo \emph{simple} (no
dirigido) cuenta caminos; el an\'alogo dirigido, que el operador no-retroceso exige, faltaba.

\begin{lemma}[Las potencias cuentan caminos]
\label{lem:walks}
Para una matriz $M$ sobre un semianillo conmutativo, la entrada $(i,j)$ de $M^k$ es la suma, sobre
los caminos de longitud $k$ de v\'ertices de $i$ a $j$, del producto de los pesos de las aristas
(\lean{pow\_apply\_eq\_sum}). Para la matriz de adyacencia $0$--$1$ de una relaci\'on decidible---un
grafo dirigido finito---esto es un conteo genuino: $\bigl((\text{adyacencia})^k\bigr)_{ij}$ es el
n\'umero de caminos dirigidos de longitud $k$ de $i$ a $j$ (\lean{relMatrix\_pow\_apply}). Ambas sin
\lean{sorry}.
\end{lemma}

La traza cuenta entonces caminos \emph{cerrados}, y---sumando la diagonal---caminos cerrados con punto
base en un v\'ertice. El mismo enunciado, reorganizado sobre el \'indice c\'iclico $\Z/(m{+}1)$, se
vuelve una suma sobre caminos \emph{c\'iclicos} (\lean{trace\_pow\_eq\_sum\_cyclicProd}), y el peso de
un camino c\'iclico es invariante bajo rotar su punto base
(\lean{cyclicProd\_comp\_finRotate}, y bajo todo el grupo de rotaci\'on). Esa invariancia por
rotaci\'on es la puerta a las geod\'esicas---caminos cerrados contados salvo rotaci\'on
c\'iclica---aunque, como la Secci\'on~\ref{sec:limits} se cuida de decir, s\'olo la puerta.

\begin{figure}[h]
\centering
\begin{tikzpicture}[scale=1.0,
  v/.style={circle,draw=ggteal,fill=ggband!50,line width=0.8pt,inner sep=1.5pt,
    minimum size=6mm,font=\small},
  e/.style={-{Stealth[length=2.2mm]},line width=1.1pt},
  ee/.style={-{Stealth[length=2.2mm]},ggamber,line width=1.6pt}]
\node[v] (a) at (90:1.35) {$0$};
\node[v] (b) at (210:1.35) {$1$};
\node[v] (c) at (330:1.35) {$2$};
\draw[ee] (a) to[bend left=12] (b);
\draw[ee] (b) to[bend left=12] (c);
\draw[ee] (c) to[bend left=12] (a);
\draw[e,ggteal!55] (b) to[bend left=12] (a);
\draw[e,ggteal!55] (c) to[bend left=12] (b);
\draw[e,ggteal!55] (a) to[bend left=12] (c);
\node[font=\small,align=center] at (0,-2.1) {un camino cerrado\\sin retroceso $w$};
\begin{scope}[xshift=4.6cm,scale=0.62]
  \foreach \k/\x in {0/0,1/2.4,2/4.8}{
    \node[v] (a\k) at ($(\x,0)+(90:1.0)$) {};
    \node[v] (b\k) at ($(\x,0)+(210:1.0)$) {};
    \node[v] (c\k) at ($(\x,0)+(330:1.0)$) {};
    \draw[ee] (a\k) to[bend left=12] (b\k);
    \draw[ee] (b\k) to[bend left=12] (c\k);
    \draw[ee] (c\k) to[bend left=12] (a\k);
    \node[font=\scriptsize] at (\x,-1.5) {$w\!\circ\!\rho^{\k}$};
  }
  \draw[{Stealth[length=2mm]}-{Stealth[length=2mm]},ggamber,line width=0.8pt]
    (0.9,1.25) to[bend left=18] (1.5,1.25);
  \draw[{Stealth[length=2mm]}-{Stealth[length=2mm]},ggamber,line width=0.8pt]
    (3.3,1.25) to[bend left=18] (3.9,1.25);
  \node[font=\scriptsize\itshape,ggamber,align=center] at (2.4,2.1)
    {una \'orbita de rotaci\'on $=$ una geod\'esica};
\end{scope}
\end{tikzpicture}
\caption{Caminos cerrados sin retroceso, y la rotaci\'on que los agrupa. \emph{Izquierda:} el
tri\'angulo $K_3$ tiene seis dardos; se muestra en \'ambar un camino cerrado sin retroceso de longitud
$3$ (los dardos inversos, en gris, retroceder\'ian). \emph{Derecha:} rotar el punto base por el
desplazamiento c\'iclico $\rho=\mathtt{finRotate}$ permuta los tres caminos con punto base de esta
\'orbita dejando fijo el peso c\'iclico (\lean{cyclicProd\_comp\_finRotate})---as\'i una sola
\emph{geod\'esica} da cuenta de una \'orbita de caminos con punto base. Aqu\'i $\tr(B^3)=6$: dos
\'orbitas (el ciclo y su inverso), tres puntos base cada una.}
\label{fig:cyclic}
\end{figure}

\paragraph{Por qu\'e es \'util, llanamente.}
``Entrada de una potencia de matriz $=$ n\'umero de caminos'' es el caballo de batalla de la teor\'ia
espectral de grafos: tiempos de mezcla, el m\'etodo de los momentos para distribuciones de
autovalores, estimaciones de expansores y las f\'ormulas de traza que este art\'iculo compone, todo
descansa en ello. Tenerlo para grafos \emph{dirigidos}, no s\'olo no dirigidos, es lo que permite que
el operador no-retroceso---un objeto inherentemente dirigido sobre aristas orientadas---se cuente
siquiera.

% ---------------------------------------------------------------
\section{Los conteos de Ihara, compuestos}
\label{sec:ihara}

\emph{Si Jacobi deriva un determinante en trazas, y Bass colapsa el determinante de Ihara a datos de
v\'ertice, ?`qu\'e dice su composici\'on?}

Dice que los conteos no-retroceso de un grafo son espectrales. Sea $B$ el operador de Hashimoto, $A$
la adyacencia y $D$ la matriz de grados; la f\'ormula de Bass~\cite{IharaBass} reza
$(1-X^2)^{|V|}\det(I-XB)=(1-X^2)^{|E|}\det(I-XA+X^2(D-I))$. Derivando ambos lados
(Teorema~\ref{thm:jacobi}, en cada determinante) y sustituyendo la identidad de Newton
(Teorema~\ref{thm:newton}) en el lado no-retroceso se a\'isla la funci\'on generatriz de conteos.

\begin{theorem}[La identidad $N_k$ de Ihara, en Lean~4]
\label{thm:ihara}
Con $Q=\det(I-XA+X^2(D-I))$, en el anillo de series de potencias,
\keyeq{\sum_k\tr(B^{k+1})X^k\;=\;\frac{\det(I-XB)\,\bigl((1-X^2)^{|V|}\bigr)'-(1-X^2)^{|E|}Q'-Q\,\bigl((1-X^2)^{|E|}\bigr)'}{(1-X^2)^{|V|}\det(I-XB)}.}
Las dos potencias de $(1-X^2)$ y ambos determinantes son unidades, as\'i que la divisi\'on es el
$\mathtt{invOfUnit}$ formal; la identidad es sin \lean{sorry} (\lean{ihara\_Nk\_explicit}). Expresa los
conteos no-retroceso $N_k=\tr(B^k)$ puramente a trav\'es de los datos de adyacencia $(A,D)$.
\end{theorem}

Y los conteos son combinatorios. Como $B$ es la matriz de adyacencia $0$--$1$ de la relaci\'on de
no-retroceso sobre dardos ($d\rightsquigarrow e$ cuando la cabeza de $d$ es la cola de $e$ y $e$ no es
el inverso de $d$), el Lema~\ref{lem:walks} se especializa:

\begin{theorem}[$\tr(B^k)$ cuenta caminos sin retroceso, en Lean~4]
\label{thm:Nk}
\keyeq{\tr(B^k)\;=\;\#\{\text{caminos cerrados sin retroceso de longitud }k,\ \text{con punto base en un dardo}\}.}
Sin \lean{sorry} (\lean{trace\_hashimoto\_pow}).
\end{theorem}

As\'i, la misma derivada que empez\'o como un hecho de una l\'inea de \'algebra lineal termina, cuatro
composiciones despu\'es, como el enunciado de que los caminos sin retroceso de longitud $k$ de un
grafo suman $\tr(B^k)$, y de que este n\'umero es computable s\'olo desde el espectro de adyacencia
del grafo.

\paragraph{Por qu\'e es \'util, llanamente.}
Los conteos y espectros no-retroceso sustentan umbrales agudos de epidemias y percolaci\'on, la
detecci\'on espectral de comunidades y---a trav\'es del an\'alogo de Ihara de la hip\'otesis de
Riemann---la definici\'on misma de un grafo de Ramanujan como expansor \'optimo. El
Teorema~\ref{thm:ihara} certifica el puente entre esos conteos y los datos de adyacencia desde los que
suelen computarse; el Teorema~\ref{thm:Nk} certifica que los conteos significan lo que se dice que
significan.

\begin{table}[h]
\centering\small
\caption{Cada ingrediente verificado por m\'aquina. Todos son sin \lean{sorry}; \lean{\#print axioms}
reporta s\'olo \lean{[propext, Classical.choice, Quot.sound]}. Fuentes:
\lean{Ihara/TraceFormula.lean} y \lean{MathlibPR/MatrixWalkCount.lean}.}
\label{tab:results}
\begin{tabular}{@{}llc@{}}
\toprule
\textbf{Nombre Lean} & \textbf{Enunciado} & \textbf{\S}\\
\midrule
\lean{Matrix.derivative\_det} & $(\det M)'=\tr(\adj M\cdot M')$ (Jacobi) & \ref{sec:jacobi}\\
\lean{one\_sub\_smul\_mul\_resolventSeries} & $(I-XM)\sum_k M^kX^k=I$ & \ref{sec:resolvent}\\
\lean{trace\_resolventSeries} & $\tr\sum_k M^kX^k=\sum_k\tr(M^k)X^k$ & \ref{sec:resolvent}\\
\lean{Matrix.charpolyRev\_logDeriv} & $(\charpolyRev M)'=-\charpolyRev M\sum_k\tr(M^{k+1})X^k$ (Newton) & \ref{sec:newton}\\
\lean{pow\_apply\_eq\_sum} & $(M^k)_{ij}=\sum_{\text{caminos}}\prod(\text{pesos})$ & \ref{sec:walks}\\
\lean{relMatrix\_pow\_apply} & conteo de caminos dirigidos, caso $0$--$1$ & \ref{sec:walks}\\
\lean{trace\_pow\_eq\_sum\_cyclicProd} & $\tr(M^{m+1})=\sum_{\text{$w$ c\'iclico}}\text{cyclicProd}$ & \ref{sec:walks}\\
\lean{cyclicProd\_comp\_finRotate\_iterate} & invariancia por rotaci\'on (grupo completo) & \ref{sec:walks}\\
\lean{ihara\_Nk\_explicit} & $\sum_k\tr(B^{k+1})X^k$ desde $(A,D)$ & \ref{sec:ihara}\\
\lean{trace\_hashimoto\_pow} & $\tr(B^k)=\#$ caminos cerrados sin retroceso & \ref{sec:ihara}\\
\bottomrule
\end{tabular}
\end{table}

\begin{table}[h]
\centering\small
\caption{Comprobaciones num\'ericas independientes (\texttt{SymPy}), igualdad simb\'olica exacta. La
fila del grafo es el tri\'angulo $K_3$: la traza de $B^k$ coincide con una enumeraci\'on por fuerza
bruta de caminos cerrados sin retroceso con punto base, $0,0,6,0,0,6$ para $k=1,\dots,6$ (el ciclo y
su inverso, tres puntos base cada uno).}
\label{tab:checks}
\begin{tabular}{@{}lll@{}}
\toprule
\textbf{Identidad} & \textbf{Instancia} & \textbf{Resultado}\\
\midrule
Jacobi $(\det M)'=\tr(\adj M\,M')$ & matriz $3\times3$ de polinomios & exacto\\
Newton (derivada log.) & $3\times3$ y $2\times2$, m\'od $X^8$ & exacto\\
$\tr(B^k)=\#\{\text{caminos c.p.\ base}\}$ & $K_3$, $k=1,\dots,6$ & $0,0,6,0,0,6$ ambos\\
Bass $(1{-}u^2)^{|V|}\det(I{-}uB)=\dots$ & $K_3$ & exacto\\
\bottomrule
\end{tabular}
\end{table}

% ---------------------------------------------------------------
\section{Lo que afirmo, y lo que no}
\label{sec:limits}

La matem\'atica es cl\'asica en todo: la f\'ormula de Jacobi, las identidades de Newton, el conteo de
caminos y el teorema de Bass son todos viejos. La contribuci\'on es la formalizaci\'on, y separo
ambas a prop\'osito.

\paragraph{Lo que afirmo como nuevo (como formalizaci\'on).}
(N1) la primera demostraci\'on verificada por m\'aquina de la que tengo noticia, en cualquier
asistente de pruebas, de la f\'ormula \emph{algebraica} de Jacobi $(\det M)'=\tr(\adj M\cdot M')$; (N2)
la forma \emph{matricial-de-trazas} de las identidades de Newton,
$(\charpolyRev M)'=-\charpolyRev M\sum_k\tr(M^{k+1})X^k$, que Mathlib ten\'ia s\'olo en el coeficiente
lineal; (N3) la serie resolvente $(I-XM)^{-1}=\sum_k M^kX^k$ y el conteo de caminos en grafos
\emph{dirigidos}, ninguno en la biblioteca. Todas son contribuciones de formalizaci\'on; ninguna es
matem\'atica nueva.

\paragraph{Y, llanamente, lo que esto no es.}
$\tr(B^k)$ cuenta caminos cerrados sin retroceso \emph{con punto base}---caminos con un dardo base
elegido---, \emph{no} geod\'esicas cerradas, que son las clases de equivalencia c\'iclica de tales
caminos, ni las geod\'esicas primas cuyo producto de Euler es la zeta de Ihara anal\'itica. La
estructura c\'iclica que las distingue \emph{s\'i} se prepara aqu\'i: la traza se reescribe como suma
sobre caminos c\'iclicos, y se demuestra que el peso es invariante bajo todo el grupo de rotaci\'on,
de modo que cada peso es funci\'on de clase sobre las \'orbitas de rotaci\'on. Pero el conteo de
\'orbitas, la contabilidad del periodo primitivo, la inversi\'on de M\"obius y el producto de Euler
que vuelven los conteos con punto base en conteos de geod\'esicas primas \emph{no} est\'an construidos.
He puesto el suelo combinatorio; la superestructura num\'erico-te\'orica est\'a trazada, no levantada.

% ---------------------------------------------------------------
\section{Implicaciones}
\label{sec:implications}

Tres clases, en orden decreciente de con cu\'anta firmeza puedo afirmarlas, y un apunte
metodol\'ogico.

\paragraph{Bloques de construcci\'on reutilizables y certificados.}
La f\'ormula algebraica de Jacobi, la serie resolvente, la identidad de Newton matricial-de-trazas y
el conteo de caminos dirigidos est\'an ahora verificados por m\'aquina y disponibles para quien
formalice polinomios caracter\'isticos, el algoritmo de Faddeev--LeVerrier, funciones zeta de grafos o
espectros de expansores. Mathlib no ten\'ia ninguno de los cuatro en forma utilizable. Es la
implicaci\'on que puedo afirmar con m\'as firmeza, y la primera a la que apuntar\'ia a un futuro
formalizador.

\paragraph{Un punto de apoyo para lo que estos desbloquean.}
La identidad de Newton es la puerta a computar coeficientes del polinomio caracter\'istico desde
trazas y de vuelta; Jacobi subyace a Faddeev--LeVerrier y a la f\'ormula de Liouville; la composici\'on
de Ihara est\'a a un paso del conteo de geod\'esicas primas que prepara. Cada uno es un hito que estas
piedras podr\'ian sostener.

\paragraph{La matem\'atica, con una nota de cautela.}
Estas identidades importan porque importan los objetos r\'io abajo---polinomios caracter\'isticos,
espectros no-retroceso, expansores de Ramanujan. Recelo de pedir prestada esa importancia: certificar
una identidad fundamental no env\'ia ning\'un algoritmo ni ninguna red. Su valor est\'a en los
cimientos certificados y en el corpus creciente de matem\'atica formalizada, no en la aplicaci\'on
inmediata; afirmo lo primero, no lo segundo.

\paragraph{Un apunte metodol\'ogico.}
Este desarrollo se llev\'o a cabo con un sistema de IA como instrumento de investigaci\'on: propon\'ia
definiciones, lemas y t\'acticas, y un lazo de compilaci\'on-como-or\'aculo verificaba o rechazaba
cada uno contra el n\'ucleo. La contabilidad cuidadosa de aqu\'i---cada afirmaci\'on dimensionada
exactamente a lo demostrado, la distinci\'on punto-base-frente-a-geod\'esica trazada con la misma
nitidez que los teoremas---es parte de lo que espero que el trabajo muestre: que la formalizaci\'on
asistida por IA puede resistir el exceso de afirmaci\'on que el m\'etodo invita, precisamente porque
el asistente de pruebas, no el asistente, tiene la \'ultima palabra.

% ---------------------------------------------------------------
\section{D\'onde esto deja de ser \'util}
\label{sec:context}

La honestidad corta por ambos lados, as\'i que enuncio los l\'imites con la misma claridad que los
resultados. El Teorema~\ref{thm:ihara} es una identidad de \emph{funciones generatrices}, no la zeta
anal\'itica; no localiza polo alguno ni nombra primo alguno. El Teorema~\ref{thm:Nk} cuenta caminos
\emph{con punto base}, y la distancia hasta las geod\'esicas---toda una teor\'ia de conteo de collares
e inversi\'on de M\"obius---es real, no cosm\'etica. La identidad de Newton se demuestra sobre un
anillo conmutativo pero no dice nada sobre cu\'ando el polinomio caracter\'istico \emph{se
descompone}; Jacobi es para matrices de polinomios, y la versi\'on suave de variedades sigue siendo un
objeto aparte. Este art\'iculo es un pu\~nado de piedras certificadas cerca del inicio de un largo
camino, y prefiero nombrar el camino antes que exagerar las piedras.

% ---------------------------------------------------------------
\section{Trabajo relacionado}
\label{sec:related}

La matem\'atica es cl\'asica: la f\'ormula de Jacobi y la de Liouville, las identidades de Newton y el
algoritmo de Faddeev--LeVerrier construido sobre ellas, la lectura de las potencias de matrices como
conteo de caminos y---para la aplicaci\'on---la funci\'on zeta de Ihara~\cite{Ihara1966}, el operador
de aristas de Hashimoto~\cite{Hashimoto1989}, la f\'ormula del determinante de Bass~\cite{Bass1992} y
el relato moderno de Terras~\cite{Terras2011}. Del lado de la formalizaci\'on, Mathlib~\cite{Mathlib2020}
provee determinantes, la adjugada, la regla de Cramer, el polinomio caracter\'istico invertido (s\'olo
en el coeficiente lineal), las identidades de Newton para polinomios sim\'etricos y conteos de caminos
para grafos simples---pero no la f\'ormula algebraica de Jacobi, ni la identidad de Newton
matricial-de-trazas, ni la serie resolvente, ni el conteo de caminos dirigidos. Este desarrollo es
compa\~nero de mis formalizaciones previas en Lean del polinomio de emparejamientos~\cite{PartI,PartII}
y de la f\'ormula de Bass~\cite{IharaBass}; no tengo noticia de una demostraci\'on previa verificada
por m\'aquina de la f\'ormula algebraica de Jacobi ni de la identidad de Newton matricial-de-trazas en
ning\'un asistente de pruebas.

% ---------------------------------------------------------------
\section{Conclusi\'on}

Part\'i del hecho m\'as elemental que pude nombrar---c\'omo se mueve un determinante cuando se mueve su
matriz---y lo segu\'i, en un asistente de pruebas, a trav\'es del diccionario de Newton entre trazas y
polinomios caracter\'isticos, y hasta los conteos de caminos cerrados de un grafo. Nada de esto es
matem\'atica nueva; la contribuci\'on es que la f\'ormula algebraica de Jacobi, la identidad de Newton
matricial-de-trazas, la serie resolvente y el conteo de caminos dirigidos est\'an ahora verificados por
m\'aquina donde no hall\'e formalizaci\'on previa, y que se componen, hasta los mismos conteos de
Ihara. El conteo de geod\'esicas primas que preparan, y la generalidad espectral plena, quedan por
delante---trazados, pero no levantados. Lo ofrezco como una perla y una invitaci\'on.

\paragraph{Artefactos.} Todas las fuentes Lean (sin \lean{sorry}, dependientes s\'olo de los tres
axiomas est\'andar) est\'an disponibles en
\begin{center}
\url{https://github.com/karlesmarin/godsil-gutman-lean}
\end{center}
en \lean{Ihara/TraceFormula.lean} (Jacobi, Newton, la composici\'on de Ihara) y
\lean{MathlibPR/MatrixWalkCount.lean} (la serie resolvente, el conteo de caminos, la estructura
c\'iclica). Verifica la huella de axiomas con, p.\,ej., \lean{\#print axioms Matrix.charpolyRev\_logDeriv}
$\rightsquigarrow$ \lean{[propext, Classical.choice, Quot.sound]}.

\begin{thebibliography}{99}
\bibitem{Bass1992} H.~Bass, \emph{The Ihara--Selberg zeta function of a tree lattice}, Internat. J.
  Math. \textbf{3} (1992), 717--797.
\bibitem{Hashimoto1989} K.~Hashimoto, \emph{Zeta functions of finite graphs and representations of
  $p$-adic groups}, en \emph{Automorphic forms and geometry of arithmetic varieties}, Adv. Stud. Pure
  Math. \textbf{15} (1989), 211--280.
\bibitem{Ihara1966} Y.~Ihara, \emph{On discrete subgroups of the two by two projective linear group
  over $\mathfrak{p}$-adic fields}, J. Math. Soc. Japan \textbf{18} (1966), 219--235.
\bibitem{Mathlib2020} The mathlib Community, \emph{The Lean mathematical library}, en \emph{Proc. 9th
  ACM SIGPLAN Int. Conf. on Certified Programs and Proofs (CPP 2020)}, 367--381.
\bibitem{PartI} C.~Mar\'in, \emph{Random Signs into Matchings: a machine-checked Godsil--Gutman
  identity in Lean~4} (2026).
\bibitem{PartII} C.~Mar\'in, \emph{Unfolding a Graph into a Tree: Heilmann--Lieb real-rootedness in
  Lean~4} (2026).
\bibitem{IharaBass} C.~Mar\'in, \emph{Folding Edges into Vertices: a machine-checked proof of Bass's
  determinant formula for the Ihara zeta in Lean~4} (2026), Zenodo DOI
  \texttt{10.5281/zenodo.20573120}.
\bibitem{Terras2011} A.~Terras, \emph{Zeta Functions of Graphs: A Stroll through the Garden},
  Cambridge Studies in Advanced Mathematics \textbf{128}, Cambridge Univ. Press, 2011.
\end{thebibliography}

\end{document}
